% \subsection{Further work}
% \textbf{Improving the methodology for the supervised case}.
% Creating better LoRAs (i.e.
% ones with either more coverage of LLM behaviour, or with more independence allowing for cleaner combinations) would allow for more robust study of the relationship between trait space and personas.
% In addition, there would be benefit to further investigation of different ways to combine, and especially to invert, the trait LoRAs.

% \textbf{Interpretability of the traits in weight space.} We would like to see exploration of where within a model different traits are shaped most, which could be done by further study into the weights themselves.
% Training LoRAs for the same trait across different models, or similar traits on the same model, could shed light into whether there is a consistent relationship between traits and weights, and whether continuity in one translates to continuity in the other.

% \textbf{Further study into the relationship between trait space and personas.} 
% Our work has primarily focused on how the behaviour of the models is impacted by our LoRAs, but little work has been done to understand the relationship between the LoRAs and the 'trait axes' themselves.
% Work exploring the relationship between amplifyer/suppressor LoRAs, or LoRAs trained on correlated traits, could shed light onto this.

% \textbf{More unsupervised exploration of personas.} From a safety perspective, understanding potential non-human traits of LLMs is useful.
% Testing alternative methods to either naturally elicit diverse personas which exist within LLMs, or to cluster and label them, would allow better study of such behaviour.
% We would also be excited to see other, entirely different methods for the unsupervised exploration of personas.